\chapter{Acute Lymphoblastic Leukemia}
\index{Acute Lymphoblastic Leukemia}

\section*{Definition and Overview}
Acute lymphoblastic leukemia (ALL) is a rapidly progressing cancer of immature lymphoid cells in the bone marrow and blood. It can occur in children or adults, with treatment and prognosis influenced by age, genetic findings, and response to therapy.

\section*{Assessment and Management}
Assessment is guided by the history, examination, and appropriate investigations. Management should be individualised by a qualified clinician and follow current Indian clinical guidance. Depending on the condition, care may include observation, medicines, procedures, surgery, rehabilitation, or specialist referral. Self-treatment should not delay evaluation of serious symptoms.

\section*{When to Seek Medical Care}
Persistent fever, pallor, bruising, bleeding, bone pain, recurrent infections, or marked fatigue should prompt urgent medical assessment.

\section*{India-specific Care}
In India, start with a qualified general physician or specialist at an appropriate clinic or hospital. Emergency symptoms should be taken to the nearest emergency department. Treatment availability varies by facility, so referral to a medical college or tertiary centre may be needed for complex disease.

\section*{Sources and Further Reading}
\begin{itemize}
    \item Indian Council of Medical Research (ICMR). Clinical and public-health guidance.
    \item Ministry of Health and Family Welfare, Government of India.
    \item World Health Organization. Relevant clinical and public-health information.
\end{itemize}

